\section{THE EXPONENTIAL AND THE FORMULA OF HAUSDORFF.}

The first research concerning the exponential map is due to E. Study and F. Engel; Engel {[}104 b{]} remarks that the exponential is not surjective for \(\mathbf{SL}(2, \mathbf{C})\) (for example \(\left( \begin{array}{cc} -1 & a \\ 0 & -1 \end{array} \right)\) is not an exponential if \(a \neq 0\) ), but that it is for \(\mathbf{GL}(n, \mathbf{C})\) , thus also for \(\mathbf{PGL}(n, \mathbf{C})\) (this latter property had already been noted by Study for \(n = 2\) ); thus \(\mathbf{SL}(2, \mathbf{C})\) and \(\mathbf{PGL}(2, \mathbf{C})\) give an example of two groups that are locally isomorphic, but which are nonetheless very different from a global point of view. Engel shows also that the exponential is surjective in the other classical groups, augmented by homotheties; this work is taken up again and continued by Maurer, Study and others, without bringing in substantial new results.

In 1899, H. Poincaré ({[}251{]}, v. III, pp.~169-172 and 173-212), tackles the study of the exponential map from a different point of view. His memoirs seem to have been hurriedly written, for in several places he affirms that every element of a connected group is an exponential, whereas he gives examples to the contrary elsewhere. His results refer mainly to the adjoint group: he shows that a semisimple element of such a group G can be the exponential of an infinity of elements of the Lie algebra \(L(G)\) , whereas a non-semisimple element may be not an exponential. If \(\operatorname{ad}(X)\) does not have an eigenvalue which is a non-zero multiple of \(2\pi i\) , then exp is étale in X. He proves also that, if U and V describe loops in \(L(G)\) , and if W is defined by continuity such that \(e^{U}.e^{V}=e^{W}\) , we do not return necessarily to the initial determination of W. He uses a formula for residues that reduces essentially to

\[
\Phi (\operatorname{ad} X) = \frac {1}{2 \pi i} \int \frac {\Phi (\xi) d \xi}{\xi - \operatorname{ad} X}
\]

where \(\operatorname{ad}(X)\) is a semisimple element whose non-zero eigenvalues have multiplicity 1, \(\Phi\) is an entire series with a sufficiently large radius of convergence, the integral being taken over a loop enclosing the eigenvalues of ad X; he studies also what happens when X tends towards a transformation having multiple eigenvalues.

The search for expressions for \(W\) as a function of \(U\) and \(V\) in the formula \(e^U.e^V = e^W\) had already, a little before the work of Poincaré, been the subject of two memoirs of Campbell {[}46 a and b{]}. As is written a little later by Baker ``\ldots the theory of Lie suggests in an obvious way that the product \(e^Ue^V\) is of the form \(e^W\) where \(W\) is an alternating series in \(U\) and \(V\ldots\)''. The subsequent work on this subject aimed to make this assertion precise and to give an explicit formula (or a method of construction) for \(W\) (``Hausdorff formula''). After Campbell and Poincaré, Pascal, Baker {[}13{]} and Hausdorff {[}152 a{]} return to the question; each considers that the proofs of his predecessor are not convincing; the main difficulty resides in what is meant by ``alternating'': is it a case of elements of the particular Lie algebra under consideration, or of universal ``symbolic'' expressions? Neither Campbell, nor Poincaré, nor Baker express themselves clearly on this point. The memoir of Hausdorff, on the contrary, is perfectly precise; he works at first in the algebra of formal associative (non commutative) series in a finite number of indeterminates and considers U, V, W as elements of this algebra. He proves the existence of W by an argument on differential equations analogous to those of his predecessors. The same argument is used by him to prove the convergence of the series when the indeterminates in it are replaced by elements of a Lie algebra of finite dimension. As Baker had remarked, and Poincaré independently, this result can be used to give a proof of the third theorem of Lie; he clarifies the correspondence between groups and Lie algebras, for example as far as the group of commutators is concerned.

In 1947, Dynkin {[}98 a{]} takes up the question again, and obtains the explicit coefficients of the Hausdorff formula, by considering from the beginning a normed Lie algebra (of finite dimension or not, over R, C or an ultrametric field). \(^{19}\)

